\documentclass{article}
\usepackage[italian]{babel}
\usepackage[utf8]{inputenc}
\usepackage{enumerate}
\usepackage{hyperref}

\title{Regolamento Greyson Allen League}
\author{}
\date{}

\begin{document}

\maketitle
\tableofcontents
\thispagestyle{empty}
\newpage
\setcounter{page}{1}

\section{Punteggi e calendario}\label{Punteggi e calendario}
\begin{enumerate}[(a)]
\item Le partite della lega si svolgono sull'app di Dunkest.
\item Il calendario è composto da: regular season, coppa di lega, play-in, playoff, Conference League, All Star Game e March Madness; le giornate adibite a ciascuna fase di ogni competizione vengono decise irrevocabilmente prima che inizi la regular season a seconda delle esigenze dettate dal calendario NBA e dal calendario di Dunkest; a tal fine possono essere inserite delle giornate di pausa.
\item La regular season inizia contestualmente alla stagione NBA, seguendo il calendario stilato da Dunkest; la sua durata non deve essere necessariamente simmetrica rispetto al numero delle squadre. 
\item A partire dall'inizio della regular season si svolge anche la coppa di lega; le giornate multiple di 3 della regular season, fino a quando il numero di giornate mancanti è uguale a 13, concorrono, oltre che per la classifica della regular season stessa, anche per una classifica parallela, stilata con gli stessi criteri utilizzati da Dunkest; dopodiché viene stilato un tabellone con le prime otto classificate, tenendo conto delle teste di serie, a eliminazione diretta in gara unica; i quarti di finale si svolgono sfalsati temporalmente in corrispondenza degli scontri diretti valevoli anche per la regular season nelle ultime 13 giornate di regular season, mentre per le semifinali e la finale sono adibite due giornate specifiche dopo la fine della regular season.
\item Dopo la fase finale della coppa di lega si disputano i play-in, della durata di due giornate; nella prima giornata, la settima classificata della regular season scontra l'ottava e la vincitrice si classifica ai playoff come settima testa di serie; contestualmente, la nona classificata sfida la decima; la squadra vincitrice di quest'ultimo incontro sfiderà, nella giornata successiva, la squadra sconfitta nel primo incontro per aggiudicarsi l'ottava testa di serie ai playoff.
\item Dopo i play-in si disputano i playoff, con le prime sei della regular season e le due qualificate tramite i play-in; i quarti di finale sono al meglio delle tre gare, le semifinali sono al meglio delle cinque gare, la finale è al meglio delle sette gare; le semifinali e la finale vengono staccate temporalmente dallo svolgimento dell'All Star Game.
\item Nelle otto giornate adibite allo svolgimento del primo turno e delle semifinali dei playoff si svolge la Conference League; il torneo si struttura in un turno di qualificazione e un tabellone finale a otto squadre, stilato tenendo conto delle teste di serie; il turno di qualificazione prevede due serie al meglio delle tre gare, in cui la quartultima della regular season sfida l'ultima della regular season per aggiudicarsi la settima testa di serie e la terzultima della regular season sfida la penultima della regular season per aggiudicarsi l'ottava testa di serie; il tabellone finale prevede i quarti di finale al meglio delle tre gare e semifinali e finale a gara unica, in cui le prime quattro teste di serie sono assegnate alle squadre eliminate al primo turno playoff, nell'ordine di classifica della regular season, la quinta testa di serie è assegnata alla squadra eliminata al secondo turno di play-in e la sesta testa di serie è assegnata alla squadra eliminata al primo turno di play-in.
\item Tra le semifinali e la finale dei playoff si svolge l'All Star Game, articolato in una fase di qualificazione e un tabellone finale a otto squadre, stilato tenendo conto delle teste di serie; accedono direttamente alla fase finale le due finaliste dei playoff, cui sono attribuite le prime due teste di serie secondo l'ordine determinato dalla classifica della regular season, la squadra vincitrice della coppa di lega, cui è attribuita la terza testa di serie, e l'altra finalista della coppa di lega, cui è attribuita la quarta testa di serie; qualora una o entrambe le squadre provenienti dalla coppa coincidano con una delle due finaliste dei playoff, i posti così liberati vengono assegnati alle squadre non ancora qualificate seguendo l'ordine della classifica della regular season; la fase di qualificazione può durare due o quattro giornate, come viene stabilito ad inizio stagione, e coinvolge le squadre non ancora qualificate, assegnando progressivamente le restanti quattro teste di serie sulla base dei punteggi più alti ottenuti nelle singole giornate, qualificando una squadra per giornata nel caso di quattro giornate e due squadre per giornata nel caso di due giornate; il tabellone finale si svolge interamente a eliminazione diretta in gara unica.
\item Al termine dei playoff si svolge la March Madness, della durata di otto giornate, a cui partecipano tutte le squadre della lega; nella prima giornata le quattro squadre con il punteggio più basso entrano in zona a rischio, mentre le restanti accedono alla zona di sicurezza; dalla seconda alla sesta giornata avvengono due eliminazioni per turno, con la prima che riguarda sempre la squadra che realizza il punteggio peggiore tra tutte quelle ancora in gara; se tale squadra appartiene alla zona a rischio, la seconda eliminazione colpisce la seconda peggiore tra le altre squadre della stessa zona, mentre le due peggiori della zona di sicurezza retrocedono in zona a rischio; se invece la prima eliminata proviene dalla zona di sicurezza, la seconda eliminazione riguarda la peggiore tra le squadre in zona a rischio, mentre la seconda peggiore della zona di sicurezza retrocede in zona a rischio; nella settima giornata vengono eliminate le due squadre che realizzano il punteggio peggiore tra le quattro ancora in gara; nell'ultima giornata si disputa la finale tra le due squadre rimaste.
\item In ogni gara dei play-in e dei playoff, la squadra che ha avuto il record migliore in regular season riceve 20 punti di vantaggio.
\item In ogni gara giocata dopo la fine della regular season, la squadra con il record migliore in regular season vince anche in caso di pareggio.
\item Formalmente, la stagione inizia al termine dell'asta (Art. \ref{Asta}) e termina alla fine della March Madness; tra una stagione e la successiva c'è l'offseason.
\end{enumerate}

\section{Roster}\label{Roster}
\begin{enumerate}[(a)]
\item Ciascun roster è composto da 15 slot regolari, due slot per i two-way contract (Art. \ref{Two-way contract e G-League}), tre slot per la G-League (Art. \ref{Two-way contract e G-League}) ed uno slot per la disabled player exception (Art. \ref{Disabled player exception}).
\item Non è obbligatorio riempire tutti gli slot del roster; tuttavia, durante tutta la stagione, devono essere presenti nel roster almeno 5 guardie, 5 ali e 3 centri secondo il listone Dunkest; a tal fine non vengono conteggiati i giocatori su cui è stata esercitata la disabled player exception (Art. \ref{Disabled player exception}) e i giocatori che non compaiono nel listone.
\item Se, durante la stagione, una squadra non possiede più il numero di giocatori minimo per qualsiasi motivo, a tale squadra viene assegnato casualmente un free agent (Art. \ref{Free agency}) da mettere sotto contratto (Art. \ref{Contratti}) fino al termine della stagione ad 1 credito; ciò avviene a prescindere dal suo salary cap (Art. \ref{Salary cap}), tranne in caso di scambio (Art. \ref{Scambi}), il quale può essere finalizzato solo se l'assegnazione del free agent non ne viola le norme; se necessario, il proprietario dovrà liberare uno slot (Art. \ref{Tagli e rilasci}).
\end{enumerate}

\section{Contratti}\label{Contratti}
\begin{enumerate}[(a)]
\item Acquistare un giocatore significa metterlo sotto contratto; i crediti relativi al contratto di ogni giocatore presente nel roster pesano sul salary cap (Art. \ref{Salary cap}) della squadra fino alla sua scadenza.
\item Cinque giorni prima dell'inizio della regular season, per ogni giocatore del proprio roster acquistato all'ultima asta (Art. \ref{Asta}), anche se acquisito tramite scambio (Art. \ref{Scambi}), va assegnato un contratto, tenendo conto di eventuali limitazioni, scegliendo tra annuale, biennale, triennale o team option (Art. \ref{Team option}).
\item Nel caso in cui ad un giocatore acquistato all'ultima asta (Art. \ref{Asta}) venga assegnato un contratto triennale, il prezzo viene ridotto del 10\%, mentre per un contratto biennale il prezzo viene ridotto del 5\%; la riduzione si applica sull'intera durata del contratto; il prezzo risultante dalla riduzione viene eventualmente arrotondato per eccesso.
\item È possibile estendere di un'ulteriore stagione i contratti con al massimo una stagione residua, relativi a giocatori acquistati all'asta (Art. \ref{Asta}) per almeno 200 crediti, originariamente triennali e non già estesi nella stagione precedente, aumentandone da subito il costo del 25\%, eventualmente arrotondato per eccesso; tale facoltà è esercitabile durante l'offseason, fino all'inizio dell'asta (Art. \ref{Asta}), da ogni squadra oppure tra la trade deadline (Art. \ref{Scambi}) e il primo termine della free agency (Art. \ref{Free agency}) dalle squadre che occupano momentaneamente una delle prime sei posizioni della regular season.
\end{enumerate}

\section{Asta}\label{Asta}
\begin{enumerate}[(a)]
\item L'asta si svolge quando mancano più di sette giorni ma meno di un mese dall'inizio della regular season NBA; per effettuare l'asta è necessario in ogni caso attendere l'uscita del listone Dunkest.
\item Durante l'asta vengono chiamati tutti i giocatori del listone Dunkest, eccetto quelli sotto contratto, seguendo l'ordine decrescente della quotazione Dunkest, con la possibilità di procedere a ritroso solo una volta giunti alla fine del listone.
\item L'asta finisce quando tutte le squadre hanno completato il proprio roster; ogni proprietario è libero di fermarsi quando vuole purché abbia firmato i 13 giocatori obbligatori.
\item Durante l'asta per un giocatore, una squadra può offrire la cifra raggiunta senza superarla se quella è la sua potenziale offerta massima per motivi di salary cap (Art. \ref{Salary cap}); in questo caso, tale squadra ha la precedenza rispetto a tutte le altre offerte di pari valore; se più squadre hanno raggiunto alla stessa cifra la loro offerta massima, si procede per sorteggio tra di esse.
\item Durante l'asta per un giocatore, una squadra può offrire la cifra raggiunta senza superarla se l'offerta è di 1 credito; tuttavia continua ad essere applicata la precedenza per le squadre che hanno raggiunto l'offerta massima; se nessuna squadra ha la precedenza, si procede per sorteggio.
\item Durante l'asta per un giocatore, una squadra può offrire la cifra raggiunta senza superarla se si tratta di una delle soglie per l'assegnazione in G-League (Art. \ref{Two-way contract e G-League}), nei two-way contract (Art. \ref{Two-way contract e G-League}) o della team option (Art. \ref{Team option}), purché si tratti di giocatori compatibili con queste possibilità; in questi casi si procede per sorteggio, senza calcolare l'eventuale precedenza dovuta al raggiungimento dell'offerta massima.
\item Una squadra può offrire dei crediti per un giocatore anche se ha tutti gli slot pieni, a patto che dopo l'acquisto il proprietario liberi uno slot (Art. \ref{Tagli e rilasci}).
\item Se durante l'asta viene tagliato o liberato (Art. \ref{Tagli e rilasci}) un giocatore che sarebbe dovuto già essere chiamato secondo l'ordine decrescente della quotazione Dunkest, l'asta su di esso si svolge immediatamente dopo il taglio.
\end{enumerate}

\section{Free agency}\label{Free agency}
\begin{enumerate}[(a)]
\item Il mercato dei free agent è aperto dal giorno della consegna dei contratti fino all'inizio dell'ultima giornata di regular season e durante tutta la March Madness.
\item L'acquisizione dei free agent avviene mediante l'app di Dunkest, ovvero tramite asta a busta chiusa prima di ciascuna giornata per i giocatori presenti nel listone Dunkest.
\item Una squadra può completare l'acquisizione di un free agent anche se ha già il roster al completo, a patto che dopo l'acquisto il proprietario liberi uno slot (Art. \ref{Tagli e rilasci}) oppure opti per un sign and trade (Art. \ref{Scambi}).
\item Ai free agent acquistati entro la trade deadline (Art. \ref{Scambi}) viene assegnato un contratto annuale con la qualifying offer (Art. \ref{Qualifying offer e restricted free agent}), mentre ai free agent acquistati successivamente viene assegnato un contratto annuale semplice.
\item Solo le squadre che occupano momentaneamente una delle prime sei posizioni della regular season possono acquisire come free agent un giocatore su cui è stato applicato il buyout (Art. \ref{Tagli e rilasci}) nella stagione in corso.
\end{enumerate}

\section{Draft}\label{Draft}
\begin{enumerate}[(a)]
\item Non prima che manchi un mese esatto all'inizio della regular season, dopo l'uscita del listone Dunkest e prima dell'asta, eventualmente anche dilazionato su più giorni, si tiene il draft dei rookie con due giri di 15 scelte.
\item Per il primo giro, l'ordine delle prime undici scelte viene determinato dalla lottery e le restanti scelte vengono assegnate alle squadre che non partecipano alla lottery nell'ordine inverso rispetto alla classifica della regular season; tutte le scelte del secondo seguono l'ordine inverso della regular season, con l'aggiunta dell'ultima scelta che viene assegnata alla squadra vincitrice della coppa di lega.
\item Nella lottery vengono assegnati due bigliettini per la squadra arrivata ultima in regular season, tre bigliettini per le altre tre squadre rimaste escluse anche dai playin, due bigliettini per le squadre eliminate al playin e un bigliettino per le squadre eliminate al primo turno playoff, anche se la loro scelta non è più sotto il loro controllo a causa di uno scambio (Art. \ref{Scambi}); inoltre viene inserito un bigliettino per una scelta che viene attribuita alla squadra vincitrice della Conference League; l'ordine delle scelte in lottery è determinato integralmente dall'estrazione dei bigliettini.
\item Possono essere scelti anche i rookie NBA che non sono presenti nel listone Dunkest, purché scelti nell'anno corrente nel draft NBA reale.
\item Non sono considerati rookie per la lega, anche se rookie per le norme NBA, quei giocatori di cui, al momento del draft, una squadra detiene il contratto, sia se in essere sia se tagliato (Art. \ref{Tagli e rilasci}), oppure la possibilità di estendere la qualifying offer (Art. \ref{Qualifying offer e restricted free agent}).
\item I giocatori scelti al primo giro vengono messi sotto contratto al seguente numero di crediti: 25, 21, 17, 14, 12, 11, 10, 9, 8, 7, 6, 5, 4, 3, 2; le scelte al secondo giro vengono firmate per 1 credito.
\item Le scelte vengono effettuate anche se il proprietario non ha uno slot disponibile tra quelli consentiti, per cui il proprietario sarà costretto a liberare uno slot adeguato (Art. \ref{Tagli e rilasci}) oppure tagliare (Art. \ref{Tagli e rilasci}) il rookie appena scelto.
\end{enumerate}

\section{Team option}\label{Team option}
\begin{enumerate}[(a)]
\item Il contratto di tutti i rookie scelti al draft prevede la team option; oltre ad essi, possono, ma non devono, essere firmati con la team option anche i rookie acquistati all'asta, indipendentemente dal prezzo di acquisto, e i giocatori under 23, ovvero quei giocatori che avranno non più di 23 anni alla fine dell'anno solare in cui inizia la stagione, ma solo se acquistati all'asta per non più di 15 crediti.
\item Dopo il primo anno di contratto ed entro un mese dall'inizio della regular season successiva, il proprietario di un giocatore con la team option può lasciar scadere il contratto oppure esercitare la team option, prolungando la durata del contratto di due anni e raddoppiandone il prezzo per l'ultimo anno; tale raddoppio avviene non appena è terminata la seconda stagione sotto contratto.
\end{enumerate}

\section{Qualifying offer e restricted free agent}\label{Qualifying offer e restricted free agent}
\begin{enumerate}[(a)]
\item L'ultimo proprietario di un giocatore può estendere la qualifying offer su di esso nell'asta della stagione successiva alla scadenza del suo contratto purché esso non sia stato tagliato o liberato (Art. \ref{Tagli e rilasci}) prima della scadenza e purché esso rientri in una delle seguenti categorie: giocatori il cui contratto prevedeva inizialmente la team option, giocatori pagati 1 credito all'ultima asta a cui è stato assegnato un contratto annuale e free agent acquistati prima della trade deadline (Art. \ref{Scambi}).
\item Estendere la qualifiyng offer su un giocatore, quando questo viene chiamato all'asta, consiste nel proporre il prolungamento del contratto per un ultimo anno al prezzo che il contratto aveva al momento della scadenza; in tal caso, quella cifra è considerata l'offerta iniziale dell'asta e non potrà essere raggiunta dalle altre squadre ma solo superata, ignorando eventuali diritti dovuti a soglie specifiche o limiti di salary cap (Art. \ref{Salary cap}); se nessuna squadra fa un'offerta superiore, la qualifying offer viene considerata accettata mentre, se viene recapitata almeno un'offerta, si considera rifiutata, il giocatore diventa restricted free agent e i diritti su di esso vengono assegnati alla squadra che aveva esteso la qualifying offer.
\item Se un giocatore è restricted free agent, all'asta possono partecipare tutte le squadre eccetto quella che ne detiene i diritti; una volta terminata, il proprietario dei diritti ha la facoltà di eguagliare l'offerta vincente per riacquisire il giocatore allo stesso prezzo; qualora eserciti tale facoltà, la squadra vincitrice dell'asta potrà scegliere se rinunciare al giocatore oppure effettuare un'ultima offerta, superiore al prezzo raggiunto fino a quel momento e senza ulteriore possibilità di rilancio; a questo punto, il proprietario dei diritti potrà decidere se trattenere il giocatore al nuovo prezzo oppure lasciarlo definitivamente alla squadra offerente a tale cifra.
\item Se la qualifying offer non viene estesa, l'asta su tale giocatore viene svolta normalmente tra tutti i partecipanti; tuttavia, se la qualifying offer viene estesa e rifiutata, il detentore dei diritti può acquistare il giocatore solo come restricted free agent.
\item Indipendentemente da altre condizioni, ad un giocatore che accetta la qualifying offer viene assegnato un contratto obbligatoriamente annuale, mentre ad un giocatore acquistato come restricted free agent può essere assegnato solamente un contratto annuale o biennale; rimane valida la riduzione del prezzo in caso venga assegnato un contratto pluriennale.
\end{enumerate}

\section{Salary cap}\label{Salary cap}
\begin{enumerate}[(a)]
\item Dall'inizio dell'offseason fino all'inizio dell'asta, il salary cap è posto a 1000 crediti, ma la firma di un giocatore scelto al draft può avvenire a prescindere dal salary cap; tuttavia, se una squadra si ritrova sopra tale soglia, essa non può effettuare altre operazioni che ne accrescano ulteriormente il salary cap.
\item All'asta, una squadra può oltrepassare i 1000 crediti ma solo con un'offerta di 1 credito, anche tramite estensione della qualifying offer; una volta superata tale soglia, la squadra può continuare ad acquistare giocatori ma solo per non più di 1 credito per ciascuno, anche tramite estensione della qualifying offer, senza alcun limite.
\item Durante la stagione, il salary cap è posto a 1100 crediti, ma la soglia dei 1000 crediti non può essere superata acquistando un free agent per più di 1 credito, né offrendo il buyout (Art. \ref{Tagli e rilasci}) ad un giocatore; inoltre, una volta superati i 1000 crediti, l'offerta massima per ogni free agent è di 1 credito e non è possibile offrire il buyout (Art. \ref{Tagli e rilasci}) ad un giocatore; se una squadra si ritrova sopra i 1100 crediti dopo la fine dell'asta, essa non deve necessariamente rientrare entro tale soglia ma, finché ciò non avviene, non può effettuare operazioni che ne accrescano ulteriormente il salary cap.
\item In qualsiasi momento, considerando come riferimento la prima stagione a doversi concludere, il salary cap della squadra per tutte le stagioni ad essa successive non può eccedere i 1000 crediti, senza eccezioni.
\item I contratti dei giocatori con la team option non pesano sul cap della squadra per le stagioni successive fino all'eventuale momento in cui vengono esercitate.
\end{enumerate}

\section{Tagli e rilasci}\label{Tagli e rilasci}
\begin{enumerate}[(a)]
\item I giocatori possono essere tagliati in ogni momento per liberare uno slot nel roster, ma continueranno a pesare sul salary cap fino alla fine del contratto. 
\item A partire dal giorno dell'esercizio delle team option e fino al primo termine della free agency, ma non durante l'asta, il proprietario di un giocatore, anche se tagliato, che non fa ufficialmente più parte di alcun roster NBA da almeno dieci giorni può rilasciarlo gratuitamente, indipendentemente dalla presenza nel listone Dunkest.
\item A partire dal giorno dell'esercizio delle team option e fino alla trade deadline (Art. \ref{Scambi}), ma non durante l'asta, il proprietario di un giocatore, anche se tagliato, sul quale NBC Sports riporti che, per squalifica o cause giudiziarie, esso abbia la certezza o rischi fortemente di saltare un numero di giornate Dunkest, calcolate al momento dell'applicazione di questa norma, pari almeno a tre quarti del numero di giornate fissate come durata della regular season, può rilasciarlo gratuitamente; in attesa che venga stabilito il calendario, la soglia è posta al 30\% della regular season NBA.
\item In offseason, fino all'inizio dell'asta, il proprietario di un giocatore che nella precedente regular season NBA ha disputato meno della metà delle partite può rilasciarlo gratuitamente, purché su di esso non sia stata esercitata la disabled player exception (Art. \ref{Disabled player exception}) durante la stagione precedente.
\item I contratti dei giocatori con la team option firmati per 1 credito possono essere rilasciati gratuitamente fino al giorno della consegna dei contratti.
\item Dalla trade deadline (Art. \ref{Scambi}) fino all'inizio della terzultima giornata di regular season, con riferimento ad un giocatore sotto contratto per ancora più di una stagione, intendendosi per stagioni residue quelle non ancora concluse, il proprietario può accordarsi con esso per il buyout, che produce i seguenti effetti: il giocatore viene tagliato, la durata del contratto viene ridotta ad una sola stagione e il costo viene concentrato integralmente.
\item L'ultimo anno di contratto di ogni giocatore su cui è stata esercitata la team option è considerato garantito solo per metà; se tali giocatori vengono tagliati prima dell'inizio dell'asta antecedente all'ultimo anno di contratto, il raddoppio del prezzo del contratto non si applica o, se è già avvenuto, viene annullato.
\item Un giocatore su cui è stata utilizzata la disabled player exception (Art. \ref{Disabled player exception}) può essere tagliato prima dell'eventuale reintegro, facendolo tornare a pesare immediatamente sul salary cap, per cui ciò non può avvenire se vengono violate le regole sul salary cap.
\item Una squadra non può riacquisire in nessun modo un giocatore che ha tagliato finché questo pesa ancora sul suo salary cap; in tale categoria rientrano anche i giocatori che sono stati rilasciati gratuitamente.
\item Il proprietario di un giocatore tagliato il cui contratto non abbia più di una stagione residua può, a partire dall'inizio dell'offseason precedente all'ultimo anno di contratto e fino alla sua scadenza e comunque massimo una sola volta per uno stesso contratto, applicare la stretch provison, suddividendo il peso residuo in due o tre rate equivalenti, arrotondate per eccesso, e distribuendole anche sugli anni successivi; ciò non può avvenire se su quel contratto è stato applicato il buyout.
\end{enumerate}

\section{Two-way contract e G-League}\label{Two-way contract e G-League}
\begin{enumerate}[(a)]
\item Tutti i giocatori con la team option sotto contratto per non più di 5 crediti possono essere inseriti nei tre slot per la G-League.
\item Tutti i giocatori il cui contratto prevede la team option, tutti i giocatori firmati come free agent nella stagione in corso e tutti i giocatori che hanno concluso la stagione precedente in G-League, purché in ciascuno dei suddetti casi risultino sotto contratto per non più di 10 crediti, possono essere inseriti nei due slot per i two-way contract.
\item Una squadra che mette sotto contratto un giocatore che rientri nei parametri per un two-way contract deve comunicare se vuole o meno inserire tale giocatore in quegli slot entro le seguenti scadenze: in ogni momento entro il giorno della consegna dei contratti per i giocatori con la team option, subito dopo l'esercizio della team option per i giocatori che hanno concluso la stagione in G-League e subito dopo l'acquisto per i free agent.
\item Una squadra che mette sotto contratto un giocatore che rientri nei parametri per la G-League deve comunicare se vuole o meno inserire tale giocatore in quegli slot in ogni momento entro il giorno della consegna dei contratti.
\item Durante l'asta, è possibile inserire un giocatore appena acquistato negli slot destinati ai two-way contract o alla G-League, purché questo rispetti tutti i requisiti; al momento della consegna dei contratti, ai giocatori che sono stati definitivamente inseriti in quegli slot deve essere assegnato un contratto con la team option.
\item Il proprietario di un giocatore in uno slot per i two-way contract può decidere in ogni momento di spostarlo tra i 15 giocatori del roster, senza modificare il tipo di contratto.
\item Il proprietario di un giocatore in uno slot per la G-League può decidere in ogni momento di spostarlo tra i 15 giocatori del roster o nei due slot per i two-way contract, senza modificare il tipo di contratto.
\item In offseason, un giocatore con la team option che ha terminato la stagione in G-League o nei two-way contract può rimanere in quegli stessi slot fino a quando non viene esercitata la team option; quando ciò avviene, il giocatore viene immediatamente promosso in uno slot superiore, per cui la team option non può essere esercitata se non vi sono slot adeguati disponibili in squadra.
\item Non è mai consentito, neanche tramite uno scambio (Art. \ref{Scambi}), lo spostamento di un giocatore da uno slot regolare ad uno slot per i two-way contract o alla G-League o da uno slot per i two-way contract alla G-League, tranne che dal giorno dell'esercizio delle team option fino alla consegna dei contratti; per i giocatori al secondo anno di contratto non vale neanche quest'ultima finestra.
\item Un giocatore al secondo anno che termina la stagione negli slot per i two-way contract viene promosso nei 15 slot regolari all'inizio dell'offseason, per cui il proprietario deve garantire uno slot disponibile o alternativamente tagliare il giocatore stesso. 
\item Se un giocatore in G-League viene utilizzato attivamente in una qualsiasi gara dei playoff, la sua team option viene esercitata immediatamente e l'anno successivo dovrà essere promosso obbligatoriamente nei 15 slot regolari all'inizio dell'offseason, per cui il proprietario a fine stagione dovrà garantire uno slot disponibile o alternativamente tagliare il giocatore stesso. 
\end{enumerate}

\section{Disabled player exception}\label{Disabled player exception}
\begin{enumerate}[(a)]
\item A partire dall'inizio dell'offseason e fino alla trade deadline, tranne che durante l'asta, il proprietario di un giocatore sul quale NBC Sports riporti che esso sia certamente indisponibile per infortunio per un numero di giornate Dunkest, calcolate al momento dell'applicazione di questa norma, pari almeno a tre quarti del numero di giornate fissate come durata della regular season, può esercitare la disabled player exception su di esso; in attesa che venga stabilito il calendario, la soglia è posta al 30\% della regular season NBA.
\item Ogni squadra dispone di uno slot extra in cui inserire un giocatore su cui è stata esercitata la disabled player exception; fino all'eventuale momento in cui verrà reintegrato nei 15 slot regolari, esso non potrà essere utilizzato e il suo contratto non peserà sul salary cap della prima stagione a doversi concludere, mentre il suo eventuale impatto sul salary cap degli anni successivi non viene alterato.
\item Al termine della stagione, un giocatore su cui è stata esercitata la disabled player exception viene reintegrato nei 15 slot regolari, a meno che il contratto non scada naturalmente; se non vi sono slot disponibili, il proprietario deve liberare uno slot o tagliare il giocatore da reintegrare.
\item Il reintegro nei 15 slot regolari di un giocatore su cui è stata esercitata la disabled player exception può, ma non deve, essere anticipato dal proprietario qualora esso torni a giocare una partita NBA prima del primo termine della free agency, purché vengano rispettate le regole sul salary cap; la decisione deve essere presa entro dieci giorni dal ritorno in campo del giocatore e comunque non oltre il primo termine della free agency.
\item I giocatori su cui è stata esercitata la disabled player exception possono normalmente essere utilizzati durante la March Madness, senza che avvenga un reintegro formale negli slot adeguati del roster.
\end{enumerate}

\section{Scambi}\label{Scambi}
\begin{enumerate}[(a)]
\item Il mercato degli scambi viene aperto al termine della stagione e viene chiuso dopo circa tre quarti delle giornate di regular season; la data esatta della trade deadline viene fissata prima dell'inizio della regular season, mentre l'ora è fissata sempre alla mezzanotte.
\item Non è possibile concludere gli scambi durante l'asta; durante la stagione, a giornata in corso, gli scambi possono essere ufficializzati ma non immediatamente applicati.
\item Insieme ai giocatori vengono scambiati i loro contratti, senza che questi subiscano variazioni; se un giocatore acquistato all'asta viene scambiato prima della consegna del contratto, è la nuova squadra ad assegnarglielo, pur rimanendo sottoposta alle stesse eventuali limitazioni.
\item Nel caso in cui venga scambiato un giocatore con un two-way contract, la nuova squadra può decidere se inserirlo nei two-way contract o tra i normali 15 giocatori del roster; nel caso in cui venga scambiato un giocatore dalla G-League, la nuova squadra può decidere se inserirlo in G-League, nei two-way contract o tra i normali 15 giocatori del roster.
\item Si possono scambiare le scelte per i successivi tre draft.
\item In offseason, si possono scambiare i diritti per estendere la qualifying offer su un giocatore.
\item Non si possono scambiare i contratti dei giocatori tagliati.
\item Si possono scambiare i giocatori su cui è stata esercitata la disabled player exception; tali giocatori dovranno essere inseriti nell'apposito slot anche nella nuova squadra.
\item Fino alla chiusura degli scambi, una squadra che firma un free agent ma non dispone immediatamente di uno slot libero può optare per una sign and trade, ossia firma del giocatore e contesuale scambio, che coinvolga o meno il giocatore appena firmato, senza dover necessariamente liberare lo slot, purché alla fine dell'operazione tutti i vincoli sul numero di slot vengano rispettati.
\end{enumerate}

\section{Protezioni e diritti di swap}\label{Protezioni e diritti di swap}
\begin{enumerate}[(a)]
\item Quando una scelta futura viene scambiata, è possibile proteggerla decidendo un numero X e indicando come contropartita una scelta in possesso della squadra che vuole inserire la protezione; la contropartita deve essere una scelta che con certezza, per giro o per annata, si manifesterà successivamente alla scelta protetta; se la scelta protetta cade nelle prime X posizioni, essa viene trasmessa al proprietario della scelta indicata come contropartita in cambio di quest'ultima.
\item È possibile, per mezzo di uno scambio, inserire il diritto di swap su una scelta futura in proprio possesso, indicando come contropartita una scelta di un'altra squadra relativa allo stesso giro dello stesso draft; se la scelta indicata come contropartita si rivela migliore della scelta su cui è stato inserito il diritto di swap, allora i proprietari invertono il possesso delle due scelte; se invece la scelta con il diritto di swap si rivela migliore della scelta indicata come contropartita, i proprietari ne mantengono normalmente il possesso.
\item Se una scelta protetta, abbinata al diritto di swap o indicata come contropartita per una protezione o per un diritto di swap viene scambiata, essa mantiene le condizioni indicate durante lo scambio precedente; diritti e doveri conseguenti a tali condizioni vengono trasferiti al nuovo possessore.
\item Una scelta protetta, abbinata al diritto di swap o indicata come contropartita per una protezione o per un diritto di swap non può essere coinvolta in altri accordi di protezione o di diritto di swap; l'unica operazione concessa è quella di eliminare o modificare l'accordo su di essa tramite uno scambio successivo tra le due squadre interessate.
\item Si possono coinvolgere negli accordi di protezione e di swap solo le scelte per i successivi tre draft, ad eccezione della scelta indicata come contropartita per una protezione, che può appartenere anche al quarto draft successivo.
\end{enumerate}

\section{Quota e montepremi}\label{Quota e montepremi}
\begin{enumerate}[(a)]
\item La quota è di 20 euro a squadra.
\item Il montepremi è destinato interamente al vincitore dei playoff, ad eccezione di una quota per ognuna delle seguenti squadre, eventualmente cumulative: il vincitore della regular season, il vincitore della coppa di lega, lo sconfitto in finale dei playoff, il vincitore dell'All Star Game e il vincitore della March Madness.
\end{enumerate}

\end{document}